% Part: quantified-modal-logic
% Chapter: syntax-and-semantics
% Section: classical-free-intro

\documentclass[../../../include/open-logic-section]{subfiles}

\begin{document}

\olfileid{qml}{syn}{csi}

\olsection{Classical and Free Semantics}

We distinguish \emph{classical} vs. \emph{free} quantified
modal logic (QML). Classical QML combines modal logic with 
classical quantification, that is, quantification as 
in first-order logic. Free QML combines modal logic with 
free quantification, that is, quantification as in free logic.

Semantics for classical QML uses \emph{fixed-domain models},
also called \emph{constant-domain models}. 
In these models, there is a single domain of objects that 
quantifiers range on. These models assume necessary existence: 
what there is is fixed across worlds.

Semantics for free QML uses \emph{variable domain models}.
In these models, quantifiers may range over different objects
at different worlds. These models assume contingent existence: 
what there is varies across worlds.

\end{document}